%!TEX root = ../main.tex

\chapter{Εφαρμογές και Μεθοδολογία}   

\section{Πειραματικές Εφαρμογές}

Η αξιολόγηση προηγμένων συστημάτων παράλληλου προγραμματισμού, τα οποία βασίζονται σε μοντέλα εργασιών (task-based paralel programming) όπως τα \textbf{StarPUPy} και \textbf{torcpy}, προϋποθέτει τη χρήση κατάλληλων υπολογιστικών σεναρίων ελέγχου (benchmarks). Τα σενάρια αυτά δεν πρέπει να είναι απλοί, ακαδημαϊκοί μικρο-κώδικες, αλλά πραγματικές, σύνθετες επιστημονικές εφαρμογές που εγείρουν διαφορετικές προκλήσεις σε επίπεδο διαχείρισης μνήμης, υπολογιστικού φόρτου και συγχρονισμού.

Για την σύγκριση των δύο προγραμματιστικών μοντέλων, χρησιμοποιήθηκαν τρεις ξεχωριστές πειραματικές εφαρμογές. Η πρώτη εφαρμογή είναι η Στοχαστική Βελτιστοποίηση με τη μέθοδο CMA-ES (Covariance Matrix Adaptation Evolution Strategy). Η δεύτερη εφαρμογή είναι Επεξεργασία Εικόνας και Ανάλυση Υφής με Συστοιχίες Φίλτρων Gabor. Τέλος, η τρίτη εφαρμογή είναι η Αποτίμηση Χρηματοοικονομικών Παραγώγων με τη μέθοδο Monte Carlo.

Στην συνέχεια, παρουσιάζεται συνοπτικά το θεωριτικό υπόβαθρο των εφαρμογών αυτών, καθώς και οι προκλήσεις που εγείρουν σε επίπεδο διαχείρισης μνήμης, υπολογιστικού φόρτου και συγχρονισμού. Έτσι ώστε να καταστεί σαφές γιατί η χρήση αυτών των εφαρμογών είναι κατάλληλη για την αξιολόγηση των μοντέλων παράλληλου προγραμματισμού που εξετάζονται.

\subsection{Στοχαστική Βελτιστοποίηση CMA-ES με Επίλυση Εξίσωσης Poisson (app00)}

\subsubsection{Αλγοριθμικό Πλαίσιο CMA-ES}

Ο αλγόριθμος CMA-ES (Covariance Matrix Adaptation Evolution Strategy) ανήκει στην κατηγορία των εξελικτικών στρατηγικών (Evolution Strategies) και αποτελεί μέθοδο βελτιστοποίησης χωρίς παραγώγους, κατάλληλη για μη-κυρτά, μη-λεία και θορυβώδη τοπία βελτιστοποίησης όπου οι κλασικές μέθοδοι gradient-based αδυνατούν. Σε κάθε γενιά $g$ παράγεται πληθυσμός $\lambda$ υποψήφιων λύσεων μέσω δειγματοληψίας από πολυμεταβλητή κανονική κατανομή \cite{HansenGECCO2014,Hansen2016,HansenCopenhagen,WikiCMAES}:

\begin{equation}
  \mathbf{x}_i^{(g+1)} = \mathbf{m}^{(g)} + \sigma^{(g)}\,\mathbf{y}_i,
  \qquad \mathbf{y}_i \sim \mathcal{N}\!\left(\mathbf{0},\,C^{(g)}\right),
  \quad i = 1,\ldots,\lambda
\end{equation}

όπου $\mathbf{m}^{(g)}\in\mathbb{R}^n$ είναι η τρέχουσα βέλτιστη εκτίμηση, $\sigma^{(g)}\in\mathbb{R}_+$ το μέγεθος βήματος (step-size) και $C^{(g)}\in\mathbb{R}^{n\times n}$ ο πίνακας συνδιακύμανσης \cite{HansenGECCO2014,WikiCMAES}.

Ο πίνακας $C$ προσεγγίζει βαθμιαία τον αντίστροφο Εσσιανό της αντικειμενικής συνάρτησης, συνδυάζοντας δύο μηχανισμούς ανανέωσης. Το \textbf{Rank-$\mu$ update} αξιοποιεί πληροφορία χωρικής κατανομής από τον επιλεγμένο πληθυσμό της τρέχουσας γενιάς, ενώ το \textbf{Rank-1 update} συσσωρεύει ιστορική κατευθυντήρια πληροφορίαμέσω ενός διανύσματος «μονοπατιού εξέλιξης» (evolutionpath) \cite{HansenPDF2016,HansenGECCO2014}.

Η τελική εξίσωση ανανέωσης του πίνακα $C$ είναι \cite{HansenGECCO2014}:

\begin{equation}
  C^{(g+1)} =
    \left(1 - c_1 - c_\mu\right)C^{(g)}
    + c_1\,\mathbf{p}_c^{(g+1)}\!\left(\mathbf{p}_c^{(g+1)}\right)^{\!\top}
    + c_\mu \sum_{i=1}^{\mu} w_i\,
      \mathbf{y}_{i:\lambda}^{(g+1)}\!\left(\mathbf{y}_{i:\lambda}^{(g+1)}\right)^{\!\top}
\end{equation}

όπου $c_1$ και $c_\mu$ είναι οι αντίστοιχοι συντελεστές μάθησης. Το μέγεθος βήματος $\sigma$ ελέγχεται ξεχωριστά μέσω της μεθόδου Cumulative Step-Size Adaptation (CSA). Η CMA-ES λειτουργεί αυστηρά βάσει ταξινόμησης (rank-based), προσδίδοντάςτης ανοχή σε μονότονους μετασχηματισμούς της αντικειμενικής συνάρτησης \cite{HansenGECCO2014,WikiCMAES}.

\subsubsection{Αντικειμενική Συνάρτηση — Εξίσωση Poisson}

Κάθε άτομο (διάνυσμα παραμέτρων $\boldsymbol{\theta}$) αξιολογείται επιλύοντας το αντίστροφο πρόβλημα της ελλειπτικής εξίσωσης Poisson \cite{appcode}:

\begin{equation}
  -\nabla^2 u(x,y) = f(x,y)
\end{equation}

Το δεξί μέλος $f$ κατασκευάζεται ως σταθμισμένο άθροισμα γκαουσιανών βάσεων $f(x,y) = \sum_k \theta_k \exp\!\left(-\|\mathbf{p} - \mathbf{s}_k\|^2 / 2\sigma_s^2\right)$, με κέντρα $\mathbf{s}_k$ σε προκαθορισμένες πηγές και συντελεστές $\boldsymbol{\theta}$. Η επίλυση γίνεται με επαναληπτική μέθοδο Jacobi σε κανονικό πλέγμα ($128\times128$ κόμβοι). Η τιμή της αντικειμενικής συνάρτησης ορίζεται ως το MSE (Mean Squared Error) μεταξύ των τιμών της υπολογισμένης λύσης $u$ σε $m=40$ σημεία-αισθητήρες και των αντίστοιχων τιμών αναφοράς, με πρόσθετο όρο κανονικοποίησης L2 \cite{appcode}.

\subsubsection{Παράλληλη Δομή και Χαρακτηρισμός}

Εντός κάθε γενιάς, οι αξιολογήσεις των $\lambda$ ατόμων είναι πλήρως ανεξάρτητες (embarra-\linebreak ssingly parallel). Ωστόσο, υφίσταται αυστηρό σημείο φράγματος συγχρονισμού (barrier synchronization) στο τέλος κάθε γενιάς: η γενιά $g+1$ δεν μπορεί να ξεκινήσει πριν ολοκληρωθούν όλες οι αξιολογήσεις της γενιάς $g$ και ανανεωθεί ο πίνακας $C$. 

Στην υλοποίηση με StarPU χρησιμοποιείται ασύγχρονη υποβολή εργασιών μέσω \texttt{starpu.task\_submit()} και \texttt{asyncio}, ενώ στην υλοποίηση με torcpy χρησιμοποιείται \texttt{torc.map()}. Αυτό το προφίλ κατατάσσει την εφαρμογή ως \textbf{synchronization-bound workload}, ιδανικό για την αξιολόγηση της εξισορρόπησης φορτίου (load balancing) και του κόστους φράγματος \cite{appcode}.


\subsection{Ανάλυση Υφής Εικόνων με Συστοιχίες Φίλτρων Gabor (app01)}

\subsubsection{Μαθηματική Διατύπωση Φίλτρου Gabor}

Τα φίλτρα Gabor είναι ζωνοπερατά φίλτρα (band-pass filters) τα οποία συνδυάζουν μια γκαουσιανή περιβάλλουσα με ένα ημιτονοειδές επίπεδο κύμα, και έχει δειχθεί ότι προσομοιώνουν την απόκριση των απλών νευρικών κυττάρων του οπτικού φλοιού . Είναι ευρύτατα χρησιμοποιούμενα στην αναγνώριση υφής (texture analysis), την ανίχνευση ακμών (edge detection) και τη διαίρεση εικόνων (image segmentation) \cite{WikiGabor,GaborMDPI}.

Για εφαρμογή υπό γωνία $\theta$, το σύστημα συντεταγμένων περιστρέφεται ως \cite{WikiGabor,Baeldung}:

\begin{equation}
  x' = x\cos\theta + y\sin\theta, \qquad
  y' = -x\sin\theta + y\cos\theta
\end{equation}

και το πραγματικό μέρος του φίλτρου (κατάλληλο για ανάλυση υφής)
εκφράζεται ως \cite{Baeldung}:

\begin{equation}
  g(x,y) = \exp\!\left(-\frac{x'^2 + \gamma^2 y'^2}{2\sigma^2}\right)
            \cos\!\left(\frac{2\pi x'}{\lambda} + \psi\right)
\end{equation}

Οι πέντε παράμετροι (μήκος κύματος $\lambda$, προσανατολισμός $\theta$, τυπική απόκλιση $\sigma$, λόγος πλευρών $\gamma$, μετατόπιση φάσης $\psi$) καθορίζουν πλήρως τη συχνότητα και κατεύθυνση χαρακτηριστικών που ανιχνεύει το φίλτρο. Κατά τη διακριτοποίηση εφαρμόζεται κανονικοποίηση μέσης τιμής (zero-mean normalization) $k = g(x,y) - \mathrm{mean}(g(x,y))$ για την εξάλειψη DC bias \cite{appcode,WikiGabor,GaborMDPI}.

\subsubsection{Συστοιχία Φίλτρων και FFT Convolution}

Κατασκευάζεται συστοιχία $n_\mathrm{scales} \times n_\mathrm{orients}$ πυρήνων Gabor (ενδεικτικά: $5 \times 8 = 40$ πυρήνες), με εκθετική επιλογή κλιμάκων $\lambda = 2^{s+1}$ και ομοιόμορφο διαμερισμό προσανατολισμών $\theta_o = o\pi/n_\mathrm{orients}$ στο $[0,\pi)$. Η τυπική απόκλιση κλιμακώνεται γραμμικά ως $\sigma = 0.56\lambda$ για σταθερό εύρος ζώνης μεταξύ κλιμάκων \cite{GaborMDPI,appcode}.

Η συνέλιξη εφαρμόζεται μέσω του Θεωρήματος Συνέλιξης στο πεδίο Fourier: με κατάλληλο zero-padding στις διαστάσεις $H = H_i + H_k - 1$, $W = W_i + W_k - 1$, η χωρική συνέλιξη $\mathcal{O}(H_i W_i H_k W_k)$ μετατρέπεται σε κατά σημείο πολλαπλασιασμό φασμάτων $\mathcal{O}(N\log N)$ \cite{appcode, WikiGabor}:

\begin{equation}
  R_\mathrm{full} = \mathrm{Real}\!\left(\mathcal{F}^{-1}(F_I \odot F_K)\right)
\end{equation}

και το αποτέλεσμα περικόπτεται στις αρχικές διαστάσεις της εικόνας.

Η ενέργεια Gabor υπολογίζεται ως \cite{appcode}:

\begin{equation}
  E(x,y) = \sqrt{\sum_{s=1}^{n_\mathrm{scales}}
                  \sum_{o=1}^{n_\mathrm{orients}}
                  R_{s,o}(x,y)^2 + 10^{-12}}
\end{equation}

Ακολουθεί λογαριθμική μετατροπή $\ln(1+E)$ για συμπίεση δυναμικού εύρους, δημιουργία ιστογράμματος $n_\mathrm{bins}$ κάδων και εξαγωγή βαθμωτού χαρακτηριστικού ως σταθμισμένο άθροισμα \cite{appcode}:

\begin{equation}
  \text{Feature Value} = \sum_{i=1}^{n_\mathrm{bins}} h[i-1] \times i
\end{equation}

\subsubsection{Παράλληλη Δομή και Χαρακτηρισμός}

Κάθε εικόνα αντιμετωπίζεται ως ανεξάρτητη εργασία, επομένως ο παραλληλισμός είναι επίπεδος (embarrassingly parallel) σε επίπεδο εικόνας χωρίς σημεία φράγματος. Το κυρίαρχο χαρακτηριστικό είναι η υψηλή πίεση στο υποσύστημα μνήμης: οι 2D FFT καταναλώνουν έντονα την κρυφή μνήμη (cache) και τον δίαυλο μνήμης (memory bus), ιδιαίτερα σε αρχιτεκτονικές NUMA. Αυτό κατατάσσει την εφαρμογή ως \textbf{memory bandwidth-bound workload}, ιδανικό για αξιολόγηση τοποθέτησης δεδομένων (data placement) και NUMA-aware scheduling \cite{appcode}.

\subsection{Αποτίμηση Χρηματοοικονομικών Παραγώγων με Monte Carlo (app02)}

\subsubsection{Στοχαστικό Μοντέλο — Γεωμετρική Κίνηση Brown}

Η εφαρμογή μοντελοποιεί την αποτίμηση ενός Δικαιώματος Προαίρεσης Καλαθιού (Basket Option) με ρήτρα Φραγμού τύπου Down-and-Out. Τα σύνθετα αυτά παράγωγα δεν διαθέτουν κλειστή αναλυτική λύση, σε αντίθεση με τα απλά Vanilla European Options που περιγράφονται από τον τύπο Black-Scholes-Merton, και αποτιμώνται αποκλειστικά μέσω στοχαστικής προσομοίωσης \cite{GoddardMC,GilesMC}.

Η εξέλιξη κάθε περιουσιακού στοιχείου $S_t$ μοντελοποιείται μέσω Στοχαστικής Διαφορικής Εξίσωσης (SDE) Γεωμετρικής Κίνησης Brown (GBM) \cite{GilesMC}:

\begin{equation}
  dS_t = r\,S_t\,dt + \sigma\,S_t\,dW_t
\end{equation}

Εφαρμόζοντας το Λήμμα Ito, η αναλυτική λύση για διακριτό χρονικό βήμα $\Delta t$ είναι \cite{GilesMC}:

\begin{equation}
  S_{t+\Delta t} = S_t \cdot
    \exp\!\left(\left(r - \tfrac{1}{2}\sigma^2\right)\Delta t
                + \sigma\sqrt{\Delta t}\,Z\right),
  \qquad Z \sim \mathcal{N}(0,1)
\end{equation}

όπου $r$ το επιτόκιο χωρίς κίνδυνο, $\sigma$ η μεταβλητότητα, $\mu = (r - \frac{1}{2}\sigma^2)\Delta t$ ο «drift» με διόρθωση Ito και $sd = \sigma\sqrt{\Delta t}$ ο «diffusion» \cite{appcode}.

\subsubsection{Συσχέτιση και Παραγοντοποίηση Cholesky}

Τα $n_\mathrm{assets}$ περιουσιακά στοιχεία του καλαθιού δεν είναι ανεξάρτητα: οι στοχαστικές διαδικασίες τους είναι συσχετισμένες μέσω πίνακα συσχέτισης $C \in \mathbb{R}^{n\times n}$, ο οποίος πρέπει να είναι Συμμετρικός και Θετικά Ορισμένος (SPD). Ο πίνακας $C$ κατασκευάζεται τυχαία με ανάμιξη τυχαίας SPD δομής και ταυτοτικού πίνακα μέσω συντελεστή $\alpha=0.15$ \cite{GoddardMC,GilesMC,appcode}:

\begin{equation}
  C_\mathrm{blended} = (1-\alpha)\,I + \alpha\,C
\end{equation}

Ακολούθως παραγοντοποιείται Cholesky ως $C = LL^\top$, όπου $L$ ο κάτω τριγωνικός «ρίζα» πίνακας. Για κάθε βήμα προσομοίωσης, ανεξάρτητα τυχαία διανύσματα $Z$ μετασχηματίζονται σε συσχετισμένες διαταραχές $dW = L \cdot Z$ \cite{GoddardMC,GilesMC,appcode}.

\subsubsection{Μη-Γραμμικές Διαταραχές («Heavy» Monte Carlo)}

Αμέσως μετά κάθε βήμα GBM εφαρμόζεται μη-γραμμική διαταραχή \cite{appcode}:

\begin{equation}
  S_t \leftarrow S_t \times
    \bigl(1.0 + 0.01\sin(S_t) + 0.01\ln(1 + |S_t|)\bigr)
\end{equation}

Αυτή η διαταραχή καταργεί τη δυνατότητα macro-stepping, υποχρεώνοντας ακριβή διακριτοποίηση σε όλα τα $n_\mathrm{steps}=128$ χρονικά βήματα. Η έντονη χρήση υπερβατικών συναρτήσεων ($\exp, \sin, \ln$) αυξάνει δραστικά τις πράξεις κινητής υποδιαστολής (FLOPs), μετατρέποντας την εφαρμογή σε \textbf{CPU-bound workload} \cite{appcode}.

\subsubsection{Συνθήκες Αποπληρωμής}

Η αξία του καλαθιού στη λήξη ορίζεται ως ο αριθμητικός μέσος \cite{appcode}:

\begin{equation}
  \mathrm{Basket}_T = \frac{1}{n_\mathrm{assets}}\sum_{i=1}^{n_\mathrm{assets}} S_{i,T}
\end{equation}

Η συνθήκη φραγμού Down-and-Out ελέγχεται σε κάθε χρονικό βήμα μέσω boolean διανύσματος: αν έστω ένα στοιχείο πέσει κάτω από $B=0.60$, το μονοπάτι «βγαίνει εκτός» (knock-out) και η αποπληρωμή μηδενίζεται \cite{AtlantisMC,appcode}:

\begin{equation}
  \mathrm{alive} \leftarrow \mathrm{alive}
    \wedge \bigwedge_{i=1}^{n_\mathrm{assets}} (S_{i,t} > B)
\end{equation}

Η τελική αποπληρωμή ακολουθεί λογική European Call και αναπροεξοφλείται στον χρόνο $t=0$ \cite{AtlantisMC}:

\begin{equation}
  \mathrm{PV} = \max(\mathrm{Basket}_T - K,\; 0) \cdot e^{-rT}
\end{equation}

\subsubsection{Παράλληλη Δομή — Batching και Parallel Reduction}

Το σύνολο $N = 3.2 \times 10^6$ μονοπατιών κατανέμεται σε $n_\mathrm{tasks} = \lceil N/\mathrm{batch} \rceil$ ανεξάρτητες παρτίδες (\texttt{batch}$=50000$), κάθε μία ως ξεχωριστή εργασία. Κάθε εργασία επιστρέφει μόνο τρία συγκεντρωτικά μεγέθη ($\sum\mathrm{PV}$, $\sum\mathrm{PV}^2$, $N_\mathrm{batch}$), αποφεύγοντας τη μεταφορά τεράστιων πινάκων μέσω διαύλου μνήμης ή δικτύου MPI. Η τελική εκτίμηση τιμής και τυπικού σφάλματος (SE) προκύπτει από παράλληλη αναγωγή (parallel reduction) των μερικών εκτιμητών \cite{appcode}:

\begin{equation}
  \mathrm{Price} = \frac{\sum\mathrm{PV}}{N_\mathrm{total}},
  \qquad
  \mathrm{SE} = \sqrt{\frac{\sum\mathrm{PV}^2/N_\mathrm{total}
                             - \mathrm{Price}^2}{N_\mathrm{total}}}
\end{equation}

\section{Μεθοδολογία Αξιολόγησης}

Για κάθε μία από τις τρεις εφαρμογές που συζητήσαμε στην προηγούμενη ενότητα, υλοποιήσαμε δύο εκδόσεις: μία με το μοντέλο \textbf{StarPUPy} και μία με το μοντέλο \textbf{torcpy}.

Η αξιολόγηση των δύο μοντέλων έγινε με βάση δύο κρητίρια: (α) την απόδοση (performance) και (β) την κλιμακωσημότητα (scalability). Η απόδοση μετρήθηκε σε επίπεδο χρόνου εκτέλεσης, ενώ η κλιμακωσιμότητα μετρήθηκε σε επίπεδο επιτάχυνσης (speedup) και μέγιστο πλήθος νήματων/διεργασιών που παραμένει υψηλή η απόδοση (parallel efficiency). 

\subsection{Εξασφάλιση Ισοδυναμίας Εκδόσεων και Δυσκολίες Υλοποίησης}

Οι εκδόσεις αυτές είναι πλήρως ισοδύναμες σε επίπεδο αλγοριθμικής λογικής και αριθμητικής ακρίβειας, ώστε να διασφαλίζεται ότι οι διαφορές στην απόδοση οφείλονται αποκλειστικά στα χαρακτηριστικά των μοντέλων παράλληλου προγραμματισμoύ. Επίσης, για να ενισχυθεί αυτή η ισοδυναμία, χρησιμοποιήθηκαν οι ίδιες μαθηματικές βιβλιοθήκες και για τις οποίες είχαμε απενεργοποιήσει, μέσω μεταβλητών περιβάλλοντος οποιαδήποτε βελτιστοποίηση είχαν ήδη εσωτερικά (π.χ. \texttt{OMP\_NUM\_THREADS=1} ή \texttt{OPENBLAS\_NUM\_THREADS=1} που απενεργοποιεί τα εσωτερικά threads των μαθηματικών βιβλιοθηκών) \cite{appcode}. 

Μόνο για την εφαρμογή της CMA-ES (app00) χρησιμοποιήθηκε η βιβλιοθήκη \texttt{numba}\footnote{Η \textbf{numba} είναι ένας open-source μεταγλωττιστής Just-In-Time (JIT) που μετατρέπει τον κώδικα Python και NumPy σε γρήγορο κώδικα μηχανής.}, ή οποία είναι κοινή και στις δύο εκδόσεις, για την επιτάχυνση του βρόχου αξιολόγησης μέσω JIT compilation. Η χρήση της \texttt{numba} ήταν απαραίτητη για να καταστεί εφικτή η εκτέλεση της εφαρμογής σε λογικά χρονικά πλαίσια. Ο λόγος που υπήρχε αυτό το πρόβλημα είναι ότι η εφαρμογή της CMA-ES περιλαμβάνει έναν βρόχο που εκτελείται για πολλές γενιές, και σε κάθε γενιά αξιολογούνται $\lambda$ άτομα, όπου κάθε αξιολόγηση απαιτεί την επίλυση της εξίσωσης Poisson. Χωρίς την επιτάχυνση της \texttt{numba}, ο χρόνος εκτέλεσης θα ήταν απαγορευτικός για την ολοκλήρωση των πειραμάτων. Οι υπλογισμοί αυτοί είναι σχετικά μικροί σε κόστος σε σχέση με την δέσμευση και απελευθέρωση του GIL της Python, και επειδή οι υπολογισμοί είναι πάρα πολύ σε πλήθος κατέληγε αυτοί η πολύ συχνή δεσμευση και απελευθέρωση του GIL να είναι ο κυρίαρχος παράγοντας που καθιστούσε την εκτέλεση της εφαρμογής μη πρακτική (φαίνομενο του GIL thrashing \footnote{Το \textbf{GIL thrashing} συμβαίνει όταν ένα πολυνηματικό (multithreaded) πρόγραμμα στην Python αναγκάζει το Global Interpreter Lock (GIL) να απελευθερώνεται και να δεσμεύεται συνεχώς με υπερβολική συχνότητα. Αντί να επιταχύνει τις εργασίες που επιβαρύνουν τον επεξεργαστή (CPU-bound tasks), αυτή η καθυστέρηση (overhead) μειώνει δραματικά την απόδοση.}). Η χρήση της \texttt{numba} επέτρεψε την εκτέλεση των υπολογισμών σε κώδικα μηχανής, παρακάμπτοντας τον GIL και επιτρέποντας την αποτελεσματική εκτέλεση των βρόχων αξιολόγησης χωρίς το φαινόμενο του GIL thrashing \cite{appcode}.

Οι υπόλοιπες δύο εφαρμογές (Image Processing και Monte Carlo) δεν είχαν τόσο αυστηρούς χρονικούς περιορισμούς, και έτσι δεν απαιτήθηκε η χρήση της \texttt{numba} για την επιτάχυνση των υπολογισμών.

\subsection{Πειραματική Διαδικασία}

Η εκτέλεση των πειραμάτων έγινε σε ένα σύστημα με δύο NUMA-nodes (2 sockets) Intel® Xeon® Gold 6530 Processor, καθένα με 32 φυσικούς πυρήνες (64 threads με Hyper-Threading), συνολικά 251 GB RAM και μέγιστη συχνότητα ρολογιού 4 GHz. 

Κάθε εφαρμογή εκτελέστηκε, για κάθε μοντέλο παράλληλου προγραμματισμoύ, με διαφορετικό πλήθος νήματων/διεργασιών (1, 2, 4, 8, 16, 32, 64) και επαναλήφθηκε 5 φορές για στατιστική αξιοπιστία κρατώντας πάντα τον μικρότερο χρόνο εκτέλεσης από κάθε συνδυασμό. 

Πιο συγκεκριμένα, οι εκδόσεις των εφαρμογών για το StarPUPy εκτελέστηκαν μόνο με threads, ενώ οι εκδόσεις για το torcpy εκτελέστηκαν με συνδυασμό με διεργασίες (processes) και threads, καθώς το torcpy είναι βασιμένο πάνω στο MPI. Πάντα, όμως το συνολικό πλήθος των νήματων/διεργασιών που χρησιμοποιήθηκε ήταν το ίδιο για κάθε μοντέλο και δεν ξεπερνούσε τον αριθμό των διαθέσιμων πόρων (64 physical cores), ώστε να είναι συγκρίσιμες οι μετρήσεις.

Επιπλέον, έγιναν δύο σετ πειραμάτων για κάθε εφαρμογή: ένα με NUMA-aware scheduling και ένα με Standard (NUMA-oblivious) scheduling. Σε κάθε περίπτωση, γίνετε χρήση των κατάλληλων σημαίων (flags) και μεταβλητών περιβάλλοντος καθενός μοντέλου πριν την εκτέλεση των εφαρμογών. Στον Πίνακα \ref{tab:execution-flags} παρουσιάζονται συνοπτικά οι παράμετροι της πειραματικής διαδικασίας \cite{starpu2025handbook,starpu_dev_manual,ibm_torcpy}.

\begin{xltabular}{\textwidth}{|l|l|X|X|}
  \hline
  \textbf{Μοντέλο} & \textbf{Λειτουργία} & \textbf{Παράμετρος} & \textbf{Περιγραφή} \\
  \hline
  \multirow{6}{*}{torcpy}
  & \multirow{4}{*}{NUMA-aware}
  & \texttt{mpirun -n N} & Εκτέλεση με N MPI διεργασίες. \\
  \cline{3-4}
  & & \texttt{-{}-map-by socket} & Αντιστοίχιση MPI διεργασιών ανά NUMA κόμβο. \\
  \cline{3-4}
  & & \texttt{-{}-bind-to core} & Δέσμευση κάθε διεργασίας σε φυσικό πυρήνα. \\
  \cline{3-4}
  & & \texttt{TORCPY\_WORKERS} & Αριθμός νημάτων εργασίας ανά MPI διεργασία, ίσος με τους πυρήνες ανά NUMA κόμβο. \\
  \cline{2-4}
  & \multirow{2}{*}{Standard}
  & \texttt{mpirun -n N} & Εκτέλεση με N MPI διεργασίες. \\
  \cline{3-4}
  & & \texttt{TORCPY\_WORKERS} & Συνολικός αριθμός νημάτων εργασίας χωρίς δέσμευση σε NUMA κόμβο. \\
  \hline
  \multirow{9}{*}{StarPUPy}
  & \multirow{6}{*}{NUMA-aware}
  & \texttt{STARPU\_NCPU} & Αριθμός CPU εργατών, ίσος με τους πυρήνες ανά NUMA κόμβο. \\
  \cline{3-4}
  & & \texttt{STARPU\_USE\_NUMA=1} & Ενεργοποίηση χαρτογράφησης μνήμης ανά NUMA κόμβο μέσω hwloc. \\
  \cline{3-4}
  & & \texttt{STARPU\_SCHED=dmda} & Χρονοπρογραμματιστής με επίγνωση κόστους μεταφοράς δεδομένων. \\
  \cline{3-4}
  & & \texttt{STARPU\_WORKERS\_GET\linebreak BIND=1} & Σεβασμός CPU masks του MPI launcher για αποφυγή υπερφόρτωσης πυρήνων. \\
  \cline{3-4}
  & & \texttt{STARPU\_LIMIT\_CPU\_N\linebreak UMA\_MEM} & Όριο μνήμης ανά NUMA κόμβο σε MB (τιμή: 32000). \\
  \cline{3-4}
  & & \texttt{STARPU\_NCUDA=0}, \texttt{STARPU\_NOPENCL=0} & Απενεργοποίηση GPU επιταχυντών. \\
  \cline{2-4}
  & \multirow{3}{*}{Standard}
  & \texttt{STARPU\_NCPU} & Συνολικός αριθμός CPU εργατών χωρίς δέσμευση NUMA τοπολογίας. \\
  \cline{3-4}
  & & \texttt{STARPU\_NCUDA=0}, \texttt{STARPU\_NOPENCL=0} & Απενεργοποίηση GPU επιταχυντών. \\
  \cline{3-4}
  \hline
  \caption{Μεταβλητές περιβάλλοντος για NUMA-aware και Standard εκτέλεση των StarPUPy και torcpy.}
  \label{tab:execution-flags}
\end{xltabular}

Η εκτέλεση των δύο αυτών σετ πειραμάτων επέτρεψε την αξιολόγηση της επίδρασης του NUMA-aware scheduling στην απόδοση και κλιμακωσιμότητα των δύο μοντέλων παράλληλου προγραμματισμού, καθώς και τη σύγκριση μεταξύ τους υπό διαφορετικές συνθήκες εκτέλεσης.

